\section{先把任务定义清楚：输入、输出与张量契约}
\label{intro:20b-task-contract}

我们已经学会用向量表示一个 Token。但“一个 Token 变成一列数字”并不是最终目标：我们需要的是\textbf{给定前文，预测下一个 Token}。在发明网络结构之前，先把任务所需的数据形状说清楚。

\subsection{先不要想网络：输入与答案到底是什么？}

取一个人为构造的短序列：
\[
\text{“春天 / 花 / 开了 / 。 ”}
\]
假设 Tokenizer 把它切成四个 Token，它们的编号分别是 $[1,3,4,2]$。
在语言建模任务中，一种常用的数据配对方式是：
\[
\underbrace{[1,\ 3,\ 4]}_{\text{输入}}\quad\longrightarrow\quad
\underbrace{[3,\ 4,\ 2]}_{\text{每个位置的目标}}.
\]
第一位置看见“春天”，目标是“花”；第二位置看见“春天、花”，目标是“开了”；第三位置看见前三个 Token，目标是句号。

注意：\textbf{训练时不需要为每个位置重新运行一遍完整模型。}如果使用因果注意力，三个位置可以一次前向计算，分别产生预测，再一起求损失。后面会完整推导这一点。

\subsection{一张形状表：每一步必须对得上}

假设词表大小 $|\mathcal V|=5$，输入有 $T=3$ 个位置，
隐藏宽度 $d_{\rm model}=2$。为方便讲解，先使用这样一组尺寸：

\begin{center}
\small
\renewcommand{\arraystretch}{1.3}
\begin{tabularx}{\linewidth}{@{}p{36mm}p{29mm}X@{}}
\toprule
数据 & 形状 & 每个轴表示什么\\
\midrule
Token ID & $T=3$ & 三个整数索引\\
Embedding 表 & $5\times2$ & 五个 Token，每个两个数字\\
隐藏状态 $X$ & $3\times2$ & 三个位置，每个两个特征\\
词表投影权重 $W_{\rm out}$ & $2\times5$ & 从隐藏维度转换到词表维度\\
Logits $Z=XW_{\rm out}$ & $3\times5$ & 每位置、对五个候选 Token 的分数\\
预测概率 & $3\times5$ & 对每一行的五个候选做 Softmax\\
\bottomrule
\end{tabularx}
\end{center}

这里 $T$ 是当前输入序列长度，$|\mathcal V|$ 是词表大小，
$d_{\rm model}$ 是网络在大多数 Block 之间传递的\textbf{隐藏向量宽度}；
它们代表三种不同的数量，不能混为一谈。

\subsection{为什么输出是“词表维度”，不是一个词编号？}

模型先给每个候选一个\textbf{分数}（logit），例如某个位置有
\[
\vect z=(0,1,2,0,0).
\]
经过 Softmax 后，五个位置得到总和为 $1$ 的概率。最大 logit 对应第三个候选；
如果使用贪心解码，就选取这个候选的 ID。

将每个 Token 的候选都保留下来，模型才容易通过交叉熵比较“正确候选”
与其他候选；选择具体 Token 发生在\textbf{解码策略}阶段，而不是
Embedding 查表阶段。

\subsection{注意区分三种“输出”}

\begin{enumerate}
\item \textbf{Block 输出：}每个 Token 都是一条 $d_{\rm model}$ 维隐藏向量，
主要用于让下一层继续处理。
\item \textbf{输出头的 logits：}每个 Token 得到一条 $|\mathcal V|$ 维分数，
用于评价候选 Token。
\item \textbf{用户看到的文字：}从概率分布选出 Token ID 后，
由 Tokenizer 的解码规则还原成文本。
\end{enumerate}

\begin{warningbox}
并非所有机器学习任务都要变成多分类。回归可以直接输出连续数值，
分类也可以用不同的损失定义。本章只讨论\textbf{自回归语言建模}的常见设计。
同样，也不是“任何函数都能由一个有限的神经网络精确表示”；通用近似结论
具有网络形式、激活函数、定义域及近似精度等条件。
\end{warningbox}

\begin{conclusion}
先约定输入是什么、答案是什么、张量有多大，
才能知道下一步为何需要 FFN、Attention 和词表投影。
\textbf{架构图是为满足输入输出契约而设计的计算流程，不是术语拼图。}
\end{conclusion}
